\chapter{Хопфовский цикл в произведении 43-мерной ячейки и КМС-носителя}
\label{ch:v10-hopf-cycle-k43-kms-product}

\section{Канонические вершины произведения}

Внутри одного $K_{43}$ третья вершина цикла либо выбиралась одним из 30
способов, либо нарушала марковскую положительность. Рассмотрим уже
существующий носитель
\begin{equation}
 \mathcal K_{43}\otimes\mathcal H_{\rm KMS},
 \qquad \dim=43\cdot6=258.
 \label{eq:v10-hopf-product-carrier}
\end{equation}
Он содержит три ортогональные линии
\begin{equation}
 |0,\Omega\rangle,
 \qquad |Y,\Omega\rangle,
 \qquad |0,s\rangle,
 \label{eq:v10-hopf-product-vertices}
\end{equation}
где $|0\rangle$ --- вакуум ячейки, $|Y\rangle$ --- гиперзарядовая метка,
$|\Omega\rangle$ --- источник КМС-репера, а $|s\rangle$ --- его синглетная
линия. Каждая вершина различена существующими типами; выбора транспортной
координаты больше нет.

Пусть $E_\times:\mathbb C^3\to\mathbb C^{258}$ имеет эти три столбца.
Тогда
\begin{equation}
 E_\times^*E_\times=I_3,
 \qquad \rank(E_\times E_\times^*)=3.
 \label{eq:v10-hopf-product-isometry}
\end{equation}
Сжатие гиперзарядового проектора равно $\operatorname{diag}(0,1,0)$, а
сжатие синглетного КМС-проектора --- $\operatorname{diag}(0,0,1)$.

\section{Канонический марковский цикл}

Вложим генератор предыдущего гейта:
\begin{equation}
 L_\times=E_\times L_\circlearrowright E_\times^*.
 \label{eq:v10-hopf-product-generator}
\end{equation}
Его сжатие точно равно $L_\circlearrowright$, полный ранг равен $2$, сумма
каждого столбца равна нулю, а отрицательных внедиагональных скоростей нет.
Стационарное состояние
\begin{equation}
 p_\times=\frac13\bigl(|0,\Omega\rangle+|Y,\Omega\rangle+|0,s\rangle\bigr)
 \label{eq:v10-hopf-product-stationary-state}
\end{equation}
нормировано и удовлетворяет $L_\times p_\times=0$.

Ток и термодинамические величины сохраняются:
\begin{equation}
 J_{\rm edge}=\frac\kappa3,
 \qquad F_\circlearrowright=3\log2,
 \qquad \sigma=\kappa\log2.
 \label{eq:v10-hopf-product-current}
\end{equation}

\section{Что исправлено и что осталось}

КМС-множитель устраняет точную дихотомию предыдущего гейта: три вершины
каноничны по типам, а генератор марковски положителен. Однако сам носитель
не выбирает общую скорость $\kappa$. Преобразование
\begin{equation}
 \kappa\mapsto c\kappa,
 \qquad t\mapsto t/c
\end{equation}
по-прежнему сохраняет стационарное состояние, силу цикла и безразмерную
эволюцию. Карта имеет ранг/ядро $2/1$.

\begin{theorem}[КМС-множитель завершает каноническое вложение цикла]
Три линии $|0,\Omega\rangle$, $|Y,\Omega\rangle$, $|0,s\rangle$ задают
каноническое изометрическое вложение ориентированного цикла в 258-мерный
носитель. Вложенный генератор сохраняет вероятность, марковскую
положительность, стационарный ток и положительное производство энтропии.
\end{theorem}

\begin{nogocriterion}[Типовой множитель не является часами]
Устранение произвольного выбора третьей вершины не фиксирует абсолютную
проводимость. Пока $\kappa$ не выбрана общим родительским функционалом,
канонический цикл задаёт только форму и направление сквозного процесса.
\end{nogocriterion}

\begin{protocol}[Статус произведения]\begin{enumerate}
\item размерность произведения: \textbf{$258$};
\item канонические типизированные вершины: \textbf{$3/3$};
\item марковская положительность: \textbf{получена};
\item стационарный ток и энтропия: \textbf{сохранены};
\item архитектура: \textbf{$10/10$};
\item реестр происхождения: \textbf{$5/7$};
\item каноническое марковское вложение: \textbf{$1/1$};
\item абсолютная проводимость: \textbf{$0/1$};
\item общий родитель скорости: \textbf{$0/1$};
\item следующий гейт: \textbf{общий родитель проводимости хопфовского цикла}.
\end{enumerate}\end{protocol}

Машинный сертификат сохранён в
\path{s2t_v10_cell_birth_four_volume_hopf_cycle_k43_kms_product_embedding_gate_results.json}.